\documentclass[12pt]{article}
\usepackage[utf8]{inputenc}
\usepackage[T1]{fontenc}
\usepackage{hyperref}
\title{Sagala Language Resources for K-12 Curriculum Development}
\author{Compiled by Claude Sonnet 5.5 (claude-sonnet-5-5) for Andr\'as Kornai}
\date{}
\begin{document}
\maketitle

\section{Introduction}
Sagala (also Sagara, Kisagara; ISO 639-3 sbm, Glottocode saga1256) is a Bantu language of Tanzania, spoken mainly in Kilosa District (Morogoro Region) and Mpwapwa District (Dodoma Region). It is distinct from the Sagalla language of Kenya, despite the similar name. Wikipedia gives 79,000 speakers (1987 figure), while Joshua Project lists a population of 267,000 for the people group. The two figures are not comparable. Documentation is thin, and I found no educational materials in two rounds of searching.

\section{Primary Educational Resources}
No primer, textbook, school programme or orthography for Sagala turned up in search.

\section{Linguistic and Reference Materials}
The only descriptive material found is general. One numerals page notes that Sagala was undocumented and undescribed, so its phoneme analysis is preliminary. The classification is also inconsistent across sources (Ruvu group in Wikipedia, Zigula-Zaramo in Joshua Project).
\begin{itemize}
\item Wikipedia, Sagala language. \url{https://en.wikipedia.org/wiki/Sagala_language}
\item Wikipedia, Sagara language. \url{https://en.wikipedia.org/wiki/Sagara_language}
\item Wikipedia, Sagara people. \url{https://en.wikipedia.org/wiki/Sagara_people}
\item Numerals in Sagala, Max Planck Institute for Evolutionary Anthropology. \url{https://lingweb.eva.mpg.de/channumerals/Sagala.htm}
\item Endangered Languages Project, Sagala bibliography. \url{https://endangeredlanguages.com/lang/4786/bibliography?hl=pt-br}
\end{itemize}

\section{Religious and Community Resources}
Joshua Project has a profile of the Sagara people group in Tanzania. No Sagala scripture was found. A Global Recordings Network page for Sagalla (code tga) turned up, but that is the Kenyan language and should not be used for Sagala.
\begin{itemize}
\item Joshua Project, Sagara in Tanzania. \url{https://joshuaproject.net/people_groups/14636/TZ}
\end{itemize}

\section{Digital and Audio Resources}
No Sagala audio, video, corpus or software was found.

\section{Curriculum Development Context}
The reference description says Sagala is endangered and not known to be taught in schools. Search supports a picture of almost no published material. Any curriculum work would first need basic documentation: a word list, an orthography agreed with the community and recorded texts. Swahili remains the school language of Tanzania. Searches in Tanzania-specific and linguistic archives under the spellings Sagara and Kisagara may find more than this dossier did.
\end{document}